\section{Teorema de Equivalencia de Lax}

\begin{frame}
	\frametitle{\insertsection}

	We first introduce an abstract framework.
	Let $\left(V,\left\|\cdot\right\|\right)$ be a Banach space,
	$V_{0}\subset V$ a dense subspace of $V$.
	Let $L\colon V_{0}\subset V\to V$ be a linear operator.
	The operator $L$ is usually unbounded and can be thought of as a
	differential operator.
	Consider the initial value problem
	\begin{equation}\label{eq:heatPDE}
		\begin{cases}\displaystyle
			\diffp{u}{t}=
			Lu      &
			\text{en }\left[0,T\right]. \\
			u=u_{0} &
			\text{en }\left\{t=0\right\}.
		\end{cases}
	\end{equation}
	This problem also represents an initial-boundary value problem with
	homogeneous boundary value conditions when they are included in
	definitions of the space $V$ and the operator $L$.
	The next definition gives the meaning of a solution of the problem
	\begin{definition}
		A function $u\colon\left[0,T\right]\to V$ is a solution of the
		initial value problem (6.2.1) if and only if
		\begin{equation*}
			\forall t\in\left[0,T\right],
			u\left(t\right)\in V_{0}:
			\lim_{\Delta t\to 0}
			\left\|
			\frac{1}{\Delta t}
			\left[u\left(t+\Delta t\right)-u\left(t\right)\right]-
			Lu\left(t\right)
			\right\|=
			0\text{ y }
			u\left(0\right)=u_{0}.
		\end{equation*}
		In the above definition, the limit in (6.2.2) is understood to be
		the right limit at $t=0$ and the left limit at $t=T$.
	\end{definition}
\end{frame}

\begin{frame}
	\frametitle{\insertsection}

	\begin{definition}
		The initial value problem (6.2.1) is well-posed if for any
		$u_{0}\in V_{0}$, there is a unique solution $u=u\left(t\right)$
		and the solution depends continuously on the initial value:
		There exists a constant $c_{0}>0$ such that if $u\left(t\right)$
		and $\overline{u}\left(t\right)$ are the solutions for the
		initial values $u_{0}$, $u_{0}\in V_{0}$, then
		\begin{equation*}
			\sup_{t\in\left[0,T\right]}
			{\left\|
				u\left(t\right)-
				\overline{u}\left(t\right)
				\right\|}_{V}\leq
			c_{0}
				{\left\|
					u_{0}-
					\overline{u}_{0}
					\right\|}_{V}.
		\end{equation*}
	\end{definition}

	From now on, we assume the initial value problem (6.2.1) is well-posed.
	We denote the solution as
	\begin{equation*}
		u\left(t\right)=
		S\left(t\right)
		u_{0},\quad u_{0}\in V_{0},\quad t\in\left[0,T\right]
	\end{equation*}
	Using the linearity of the operator $L$, it is easy to see that the
	solution operator $S\left(t\right)$ is linear.
	From the continuous dependence property (6.2.3), we have
	\begin{align*}
		\sup_{t\in\left[0,T\right]}
		{\left\|
			S\left(t\right)
			\left[u_{0}-\overline{u}_{0}\right]
		\right\|}_{V} & \leq
		c_{0}
			{\left\|
				u_{0}-
				\overline{u}_{0}
		\right\|}_{V}.       \\
		\forall u_{0}\in V_{0}:
		\sup_{t\in\left[0,T\right]}
		{\left\|
			S\left(t\right)
			u_{0}
		\right\|}_{V} & \leq
		c_{0}
			{\left\|
				u_{0}-
				\right\|}_{V}.
	\end{align*}
	By Theorem 2.4.1, the operator
	$S\left(t\right)\colon V_{0}\subset V\to V$ can be uniquely
	extended to a linear continuous operator $S\left(t\right)\colon V\to V$ with
	\begin{equation*}
		\sup_{t\in\left[0,T\right]}
		{\left\|
			S\left(t\right)
			\right\|}_{V} \leq
		c_{0}.
	\end{equation*}
\end{frame}

\begin{frame}
	\frametitle{\insertsection}

	\begin{definition}
		.
	\end{definition}
\end{frame}
